\section{Discussion}

\label{sec:discussion}
% ══════════════════════════════════════════════════════════════════════════════

\subsection{Market Implications}

The 30.7\% over-crediting correction has significant market implications.
If applied across the entire VCM, approximately 200 million credits out
of the 650 million issued for NBS projects may not represent genuine
climate benefit. This supports the growing consensus that quality
assurance---not volume growth---is the primary challenge for carbon
markets \citep{badgley2022}.

However, fewer high-quality credits command premium prices. The voluntary
market is increasingly stratified, with ``premium'' credits verified by
advanced dMRV commanding 2--5x price premiums over standard credits.
This price signal can drive adoption of ZooMRV-class verification
infrastructure.

\subsection{Advantages of Digital MRV}

\begin{enumerate}
  \item \textbf{Transparency}: Open-source models and publicly
    verifiable data chains eliminate the ``black box'' problem of
    traditional verification.
  \item \textbf{Scalability}: Satellite-based monitoring scales to
    millions of hectares without proportional cost increases.
  \item \textbf{Timeliness}: Near-real-time reversal detection enables
    rapid response and market-responsive credit adjustment.
  \item \textbf{Consistency}: Algorithmic baselines and verification
    eliminate auditor variability and subjective judgment.
  \item \textbf{Cost efficiency}: After initial infrastructure
    investment, per-credit verification costs are estimated at
    \$0.02--0.05, compared to \$0.50--2.00 for traditional audits.
\end{enumerate}

\subsection{Limitations}

\begin{enumerate}
  \item \textbf{Below-canopy changes}: Satellite-based monitoring
    cannot detect forest degradation that occurs beneath the canopy
    (selective logging of understory, biodiversity loss). Acoustic and
    camera trap monitoring can complement.
  \item \textbf{Soil carbon}: CarbonNet models above-ground biomass
    only. Soil organic carbon, which represents the majority of carbon
    in peatlands and grasslands, requires separate modeling approaches.
  \item \textbf{Cloud cover}: Persistent cloud cover in tropical regions
    reduces optical satellite data availability. Sentinel-1 SAR provides
    all-weather coverage but with lower sensitivity to biomass.
  \item \textbf{GEDI data gaps}: GEDI's sampling pattern leaves gaps
    that are filled by extrapolation, introducing uncertainty in areas
    without direct LiDAR measurements.
  \item \textbf{Additionality}: Dynamic baselines address over-crediting
    but do not fully resolve the fundamental challenge of proving that
    conservation would not have occurred without carbon finance.
\end{enumerate}


% ══════════════════════════════════════════════════════════════════════════════
